\subsection{Linear Approximation}

Let $f: \mathbb{R}^n \to \mathbb{R}^m$ be defined by $f(\mathbf{x})=(f_1(\mathbf{x}), f_2(\mathbf{x}), \dots, f_m(\mathbf{x}))$, where each component $f_i: \mathbb{R}^n \to \mathbb{R}$ is a real-valued function. 

Regardless of whether $f$ is linear or non-linear, a natural question arises: can it be locally approximated by a linear function? And if so, how accurate is this approximation?

Ultimately, given a point $\mathbf{x}_0 \in \mathbb{R}^n$, our goal is to approximate $f$ near $\mathbf{x}_0$ using its linearization:
$$f(\mathbf{x}) \approx f(\mathbf{x}_0) + Df(\mathbf{x}_0) (\mathbf{x} - \mathbf{x}_0)$$


\begin{definition}[Differentiability]
Let $E \subseteq \mathbb{R}^n$, $f: E \to \mathbb{R}^m$ be a function, and $\mathbf{x}_0$ be an interior point of $E$. Let $L: \mathbb{R}^n \to \mathbb{R}^m$ be a linear transformation. We say $f$ is differentiable at $\mathbf{x}_0$ with derivative $L$ if 
$$\lim_{\mathbf{x} \to \mathbf{x}_0} \frac{\|f(\mathbf{x}) - f(\mathbf{x}_0) - L(\mathbf{x}-\mathbf{x}_0)\|}{\|\mathbf{x}-\mathbf{x}_0\|} = 0.$$
Equivalently,
$$\lim_{\mathbf{h}\to\mathbf{0}}\frac{\| f(\mathbf{x}_0+\mathbf{h})-f(\mathbf{0})-L(\mathbf{h}) \|}{\|\mathbf{h}\|}=0.$$
\label{differentiablity}
\end{definition}

Therefore, multivariable differentiability is defined such that the relative error between the function and its linear approximation vanishes as the distance approaches zero.

\begin{theorem}
Let $E \subseteq \mathbb{R}$, $f: E \to \mathbb{R}$ be a function, $L \in \mathbb{R}$, and $x_0$ be an interior point of $E$. TFAE:
\begin{enumerate}
    \item $f$ is differentiable at $x_0$ and $f'(x_0) = L$.
    \item $\lim_{x \to x_0} \frac{|f(x) - f(x_0) - L(x - x_0)|}{|x - x_0|} = 0$.
\end{enumerate}
\end{theorem}

\begin{pfbox}

For $x\neq x_0$, write
\[\frac{f(x) - f(x_0) }{x - x_0}= L + \frac{f(x) - f(x_0) - L(x - x_0)}{x - x_0}.\]
Rearrange $L$ to the left, take the absolute value and limit $x\to x_0$, and we can show the limit on the both sides imply each other.
The second expression is exactly the special case of definition \eqref{differentiablity}.

\end{pfbox}



\begin{theorem}[Uniqueness of derivatives]
Let $\mathbf{x}_0 \in E$ be an interior point of $E$. 
Suppose $L_1, L_2: \mathbb{R}^n \to \mathbb{R}^m$ are linear transformations such that $f$ is differentiable at $\mathbf{x}_0$ with derivative $L_1$ and also differentiable at $\mathbf{x}_0$ with derivative $L_2$. 
Then $L_1 = L_2$.
\end{theorem}


\begin{pfbox}

For $\mathbf{h}\ne \mathbf{0}$,
\[\frac{\| L_1(\mathbf{h}) - L_2(\mathbf{h}) \|}{\|\mathbf{h}\|}\le
\frac{\| f(\mathbf{x}_0+\mathbf{h})-f(\mathbf{0})-L_1(\mathbf{h}) \|}{\|\mathbf{h}\|} +
\frac{\| f(\mathbf{x}_0+\mathbf{h})-f(\mathbf{0})-L_2(\mathbf{h}) \|}{\|\mathbf{h}\|}.
\]

By sandwich, $\lim_{\mathbf{h}\to\mathbf{0}}\frac{\| L_1(\mathbf{h}) - L_2(\mathbf{h}) \|}{\|\mathbf{h}\|}=0.$

Let $\mathbf{h}=t\mathbf{v}, t\ne 0$, we have
\begin{align*}
\lim_{\mathbf{\mathbf{v}}\to\mathbf{0}}\frac{\| L_1(t\mathbf{v}) - L_2(t\mathbf{v}) \|}{\| t\mathbf{v}\|} = 
\lim_{\mathbf{\mathbf{v}}\to\mathbf{0}}\frac{|t| \|L_1(\mathbf{v}) - L_2(t\mathbf{v})\|}{|t| \|\mathbf{v}\|} \quad (\text{$L_1$ and $L_2$ are linear.})
\end{align*}

Hence, for $\mathbf{v}\ne\mathbf{0}$, $L_1(\mathbf{v})=L_2(\mathbf{v})$, which is also true when $\mathbf{v}=\mathbf{0}.$

\end{pfbox}

We may prove by contradiction that $f$ is not differnetiable by the above theorem.

\begin{theorem}
If $f$ is differentiable at $\mathbf{x}_0$, then $f$ is continuous at $\mathbf{x}_0$.
\end{theorem}

\begin{pfbox}

For $\mathbf{x}\neq \mathbf{x}_0$, write
\[\|f(\mathbf{x}) - f(\mathbf{x}_0)-L(\mathbf{x} - \mathbf{x}_0)\| = \underbrace{\frac{\|f(\mathbf{x}) - f(\mathbf{x}_0) - L(\mathbf{x} - \mathbf{x}_0)\|}{\|\mathbf{x} - \mathbf{x}_0\|}}_{\to 0 \text{ by assumption}} \underbrace{\|\mathbf{x} - \mathbf{x}_0\|}_{\to 0 \text{ as } \mathbf{x}\to\mathbf{x}_0}.\]
Next, use
\[\|f(\mathbf{x}) - f(\mathbf{x}_0)\|\le \underbrace{\|f(\mathbf{x}) - f(\mathbf{x}_0)-L(\mathbf{x} - \mathbf{x}_0)\|}_{\to 0 \text{ from the previous equality}} + \underbrace{\|L(\mathbf{x} - \mathbf{x}_0)\|}_{\to 0 \text{ as } \mathbf{x}\to\mathbf{x}_0}.\]
By Sandwich, $\lim_{\mathbf{x}\to\mathbf{x}_0}f(\mathbf{x})=f(\mathbf{x}_0).$

\end{pfbox}


\begin{definition}[Directional Derivative]
Let $E \subseteq \mathbb{R}^n$, $f: E \to \mathbb{R}^m$ be a function, and $\mathbf{x}_0$ be an interior point of $E$.
Fix $\mathbf{v}\in\mathbb{R}^n$, if the limit
$$D_\mathbf{v} f(\mathbf{x}_0) = \lim_{t \to 0} \frac{f(\mathbf{x}_0 + t\mathbf{v}) - f(\mathbf{x}_0)}{t}$$
exists, we call this limit the directional derivative of $f$ at $\mathbf{x}_0$ in the direction $\mathbf{v}.$
\end{definition}


\begin{definition}[Partial Derivative]
Let $E \subseteq \mathbb{R}^n$, $f: E \to \mathbb{R}^m$ be a function, and $\mathbf{x}_0$ be an interior point of $E$.
Let $1\le j\le n$, if the limit
$$\frac{\partial f}{\partial x_j}(\mathbf{x}_0) = \lim_{t \to 0} \frac{f(\mathbf{x}_0 + t\mathbf{e}_j) - f(\mathbf{x}_0)}{t}$$
exists, we call this limit the partial derivative of $f$ with respect to the variable $x_j$.
\end{definition}

So partial derivative is just a special case of directional derivative.

\begin{theorem}
If $f$ is differentiable at $\mathbf{x}_0$, then $f$ is differentiable in the direction $\mathbf{v}$ at $\mathbf{x}_0$, and:
$$D_\mathbf{v} f(\mathbf{x}_0) = f'(\mathbf{x}_0)\cdot\mathbf{v}.$$
\end{theorem}

\begin{pfbox}

For $\mathbf{h}\ne 0,$ let $\mathbf{h}=t\mathbf{v},\, t\ne 0,$ write
\[
\frac{f(\mathbf{x}_0 + t\mathbf{v})-f(\mathbf{x}_0)}{t} 
= \frac{f(\mathbf{x}_0 + t\mathbf{v})-f(\mathbf{x}_0)-L(t\mathbf{v})}{t}
+ \frac{L(t\mathbf{v})}{t}.
\]
Rearrange $L(\mathbf{v})$ to the left, and take the Euclidean norm and $t\to 0$, the RHS converge to 0 by differentiablity ($\|\mathbf{v}\|\ne 0$). So we have
\[
\lim_{t\to 0}\frac{f(\mathbf{x}_0 + t\mathbf{v})-f(\mathbf{x}_0)}{t} = L(\mathbf{v}) = f^\prime(\mathbf{x}_0)\cdot\mathbf{v}.
\]

\end{pfbox}



We may also disprove the differentiability of $f$ provided its directional derivative exists.
The next question is: how to write the linear map $f'(\mathbf{x}_0)$ explicitly?

\begin{theorem}
Let $E \subseteq \mathbb{R}^n$, $f: E \to \mathbb{R}^m$ be a function, and $\mathbf{x}_0$ be an interior point of $E$.
If all the partial derivatives of $f$ exist on $E$ and are continuous at $\mathbf{x}_0$, then $f$ is differentiable at $\mathbf{x}_0$, and
$$f'(\mathbf{x}_0)\cdot\mathbf{v} = \sum_{j=1}^n v_j \frac{\partial f}{\partial x_j}(\mathbf{x}_0).$$
\end{theorem}


\begin{pfbox}

\noindent \textbf{Step 1: Setup}\\
Let $L(\mathbf{v}) = \sum_{j=1}^n v_j \frac{\partial f}{\partial x_j}(\mathbf{x}_0)$ be our candidate for the derivative.
Since $\mathbf{x}_0$ is an interior point of $E$, there exists $r > 0$ such that the ball $B(\mathbf{x}_0, r) \subseteq E$.

For $\epsilon > 0$, since each partial derivative $\frac{\partial f_i}{\partial x_j}$ is continuous at $\mathbf{x}_0$, there exists $\delta_{i,j} > 0$ such that $\|\mathbf{x} - \mathbf{x}_0\| < \delta_{i,j}$ implies
$$\left| \frac{\partial f_i}{\partial x_j}(\mathbf{x}) - \frac{\partial f_i}{\partial x_j}(\mathbf{x}_0) \right| < \frac{\epsilon}{mn}$$
for all $1 \le i \le m$ and $1 \le j \le n$.

Let $\delta = \min\{r, \delta_{i,j} \mid 1 \le i \le m, 1\le j \le n\}$.
Now, fix $\mathbf{x} \in E$ such that $\|\mathbf{x} - \mathbf{x}_0\| < \delta$. Let $\mathbf{v} = \mathbf{x} - \mathbf{x}_0 = \sum_{j=1}^n v_j \mathbf{e}_j$. Note that $|v_j| \le \|\mathbf{x} - \mathbf{x}_0\| < \delta$.
We construct a path from $\mathbf{x}_0$ to $\mathbf{x}$ along the coordinate axes. Define the sequence of points:
\begin{align*}
\mathbf{x}^{(0)} &= \mathbf{x}_0 \\
\mathbf{x}^{(1)} &= \mathbf{x}_0 + v_1 \mathbf{e}_1 \\
\mathbf{x}^{(2)} &= \mathbf{x}_0 + v_1 \mathbf{e}_1 + v_2 \mathbf{e}_2 \\
&\vdots \\
\mathbf{x}^{(j)} &= \mathbf{x}^{(j-1)} + v_j \mathbf{e}_j \\
&\vdots \\
\mathbf{x}^{(n)} &= \mathbf{x}_0 + \sum_{j=1}^n v_j \mathbf{e}_j = \mathbf{x}
\end{align*}

We can express the difference $f(\mathbf{x}) - f(\mathbf{x}_0)$ as a telescoping sum:
$$f(\mathbf{x}) - f(\mathbf{x}_0) = \sum_{j=1}^n \left( f(\mathbf{x}^{(j)}) - f(\mathbf{x}^{(j-1)}) \right)$$

For each component $i$, consider $f_i(\mathbf{x}^{(j)}) - f_i(\mathbf{x}^{(j-1)})$. 
Let $\varphi_i^j(t) = f_i(\mathbf{x}^{(j-1)} + t v_j \mathbf{e}_j)$, then $f_i(\mathbf{x}^{(j)}) - f_i(\mathbf{x}^{(j-1)}) = \varphi^j_i(1)-\varphi^j_i(0)$.
By the Mean Value Theorem, $\exists t_j\in (0,1)$ such that $\varphi^j_i(1)-\varphi^j_i(0) = (\varphi^j_i)'(t_j)$. 

\vspace{0.8em}
\noindent \textbf{Step 2: Continuity of $f_i$ along the coordinate $\mathbf{x}_j$}\\
Applying Lemma \eqref{coordinate chain rule}: 
$$
f_i(\mathbf{x}^{(j)}) - f_i(\mathbf{x}^{(j-1)}) = (\varphi^j_i)'(t_j) = \frac{\partial f_i}{\partial x_j}(\boldsymbol{\xi}_j) v_j,
$$ 
where $\boldsymbol{\xi}_j = \mathbf{x}^{(j-1)} + t_j v_j \mathbf{e}_j$. 

Note that
\begin{align*}
\|\mathbf{x}^{(j-1)} + t_j v_j \mathbf{e}_j - \mathbf{x}_0\| &= \|v_1 \mathbf{e}_1 + v_2 \mathbf{e}_2 + \dots + v_{j-1} \mathbf{e}_{j-1} + t_j v_j \mathbf{e}_j\| \\
&\le \left\| \sum_{i=1}^n v_i \mathbf{e}_i \right\| \quad (\because |t_j| < 1) \\
&= \|\mathbf{x} - \mathbf{x}_0\| \le \delta,
\end{align*}
which implies
\begin{align*}
\left| f_i(\mathbf{x}^{(j)}) - f_i(\mathbf{x}^{(j-1)}) - \frac{\partial f_i}{\partial x_j}(\mathbf{x}_0) v_j \right| 
&= \left| \frac{\partial f_i}{\partial x_j}(\mathbf{x}^{(j-1)} + t_j v_j \mathbf{e}_j) - \frac{\partial f_i}{\partial x_j}(\mathbf{x}_0) \right| |v_j| \\
&< \frac{\epsilon}{mn} |v_j|.
\end{align*}

\vspace{0.8em}
\noindent \textbf{Step 3: Bound the error for each component $f_i$}\\
\begin{align*}
\left| f_i(\mathbf{x}) - f_i(\mathbf{x}_0) - \sum_{j=1}^n \frac{\partial f_i}{\partial x_j}(\mathbf{x}_0) v_j \right| &= \left| \sum_{j=1}^n \left( f_i(\mathbf{x}^{(j)}) - f_i(\mathbf{x}^{(j-1)}) - \frac{\partial f_i}{\partial x_j}(\mathbf{x}_0) v_j \right) \right| \\
&\le \sum_{j=1}^n \left| \frac{\partial f_i}{\partial x_j}(\boldsymbol{\xi}_j) - \frac{\partial f_i}{\partial x_j}(\mathbf{x}_0) \right| |v_j| \\
&< \sum_{j=1}^n \frac{\epsilon}{mn} |v_j| \\
&\le \sum_{j=1}^n \frac{\epsilon}{mn} \|\mathbf{x} - \mathbf{x}_0\| = \frac{\epsilon}{m} \|\mathbf{x} - \mathbf{x}_0\|
\end{align*}

\vspace{0.8em}
\noindent \textbf{Step 4: Extend this to the vector-valued function $f$}\\
\begin{align*}
\|f(\mathbf{x}) - f(\mathbf{x}_0) - L(\mathbf{x} - \mathbf{x}_0)\| &\le \sum_{i=1}^m \left| f_i(\mathbf{x}) - f_i(\mathbf{x}_0) - \sum_{j=1}^n \frac{\partial f_i}{\partial x_j}(\mathbf{x}_0) v_j \right| \\
&< \sum_{i=1}^m \frac{\epsilon}{m} \|\mathbf{x} - \mathbf{x}_0\| = \epsilon \|\mathbf{x} - \mathbf{x}_0\|
\end{align*}

Dividing both sides by $\|\mathbf{x} - \mathbf{x}_0\|$, we get:
$$\frac{\|f(\mathbf{x}) - f(\mathbf{x}_0) - L(\mathbf{x} - \mathbf{x}_0)\|}{\|\mathbf{x} - \mathbf{x}_0\|} < \epsilon$$
Since this holds for any $\epsilon > 0$ when $\|\mathbf{x} - \mathbf{x}_0\|$ is sufficiently small, we conclude:
$$\lim_{\mathbf{x} \to \mathbf{x}_0} \frac{\|f(\mathbf{x}) - f(\mathbf{x}_0) - L(\mathbf{x} - \mathbf{x}_0)\|}{\|\mathbf{x} - \mathbf{x}_0\|} = 0$$
This proves that $f$ is differentiable at $\mathbf{x}_0$, and its derivative is given by $L$.

\end{pfbox}

\begin{lemma}[Chain rule along a coordinate segment]
Let $E \subset \mathbb{R}^n$ and $f: E \to \mathbb{R}^m$. Fix $1 \le i \le m$ and $1 \le j \le n$. Define
$$ \varphi_i^j(t) := f_i(\mathbf{z} + t v \mathbf{e}_j), \quad 0 \le t \le 1, $$
where $\mathbf{z} \in E$ and $v \in \mathbb{R}$. Assume that for some $t_0 \in (0, 1)$ there exists $\rho > 0$ such that
$$ \mathbf{z} + (t_0 + h)v \mathbf{e}_j \in E \quad \text{for all } |h| < \rho, $$
and that the partial derivative $\frac{\partial f_i}{\partial x_j}$ exists at the point $\mathbf{p} := \mathbf{z} + t_0 v \mathbf{e}_j$. Then $\varphi_i^j$ is differentiable at $t_0$ and
$$ (\varphi_i^j)'(t_0) = \frac{\partial f_i}{\partial x_j}(\mathbf{p}) v. $$
\label{coordinate chain rule}
\end{lemma}


\begin{pfbox}

Let 
\begin{align*}
(\varphi_i^j)'(t_0) = \lim_{h\to 0}\frac{\varphi_i^j(t_0+h)-\varphi_i^j(t_0)}{h}
&=\lim_{h\to 0}\frac{f_i(\mathbf{z} + (t_0+h) v \mathbf{e}_j)-f_i(\mathbf{z} + t_0 v \mathbf{e}_j)}{t}\\
&=\lim_{h\to 0}\frac{f_i(\mathbf{p}+hv\mathbf{e}_j)-f_i(\mathbf{p})}{h}.
\end{align*}
If $v=0$, $f_i(\mathbf{p}+hv\mathbf{e}_j)-\mathbf{p}=0.$
Now consider $v\ne 0$. 

For arbitrarily small $h$, $\mathbf{z} + (t_0 + h)v \mathbf{e}_j=\mathbf{p}+hv\mathbf{e}_j\in B(\mathbf{p},\rho)\subseteq E.$ 
Set $hv=s$, the above quotient becomes:
\[\lim_{s\to 0}\frac{f_i(\mathbf{p}+s\mathbf{e}_j)-f_i(\mathbf{p})}{s}\cdot v = \frac{\partial f_i}{\partial x_j}(\mathbf{p})v,\]
which shows that $\varphi_i^j$ is differentiable at $t_0$ and
$(\varphi_i^j)'(t_0) = \frac{\partial f_i}{\partial x_j}(\mathbf{p}) v.$

\end{pfbox}



\begin{remark}
Let $E \subseteq \mathbb{R}^n$, $f: E \to \mathbb{R}^m$ be a function, and $\mathbf{x}_0$ be an interior point of $E$.
If $f$ is differentiable at $\mathbf{x}_0$, by defintion, there exists a linear map such that the relative error converge to zero near $\mathbf{x}_0.$
Then we call this linear map (total) derivative of $f$ at $\mathbf{x}_0,$ denoted by

\begin{align*}
Df(\mathbf{x}_0) = f^\prime(\mathbf{x}_0) &=
\begin{pmatrix}
- & \nabla f_1(\mathbf{x}_0) & - \\
- & \nabla f_2(\mathbf{x}_0) & - \\
  & \vdots & \\
- & \nabla f_m(\mathbf{x}_0) & - 
\end{pmatrix}_{m \times n}\\ \\ &= 
\begin{pmatrix}
| & | & & | \\
\frac{\partial f}{\partial x_1}(\mathbf{x}_0) & \frac{\partial f}{\partial x_2}(\mathbf{x}_0) & \cdots & \frac{\partial f}{\partial x_n}(\mathbf{x}_0) \\
| & | & & |
\end{pmatrix}_{m \times n}.    
\end{align*}
\end{remark}

Hence, we have:
\begin{enumerate}
    \item The partial derivative is simply the directional derivative along the standard basis vector $\mathbf{e}_j:$
    $$ \frac{\partial f}{\partial x_j}(\mathbf{x}_0) = D_{\mathbf{e}_j} f(\mathbf{x}_0) = f'(\mathbf{x}_0) \mathbf{e}_j $$
    
    \item If $f$ is differentiable, then for any direction vector $\mathbf{v} = \sum_{j=1}^n v_j \mathbf{e}_j$, the directional derivative can be expanded as:
    $$ D_{\mathbf{v}} f(\mathbf{x}_0) = f'(\mathbf{x}_0) \mathbf{v} = f'(\mathbf{x}_0) \left( \sum_{j=1}^n v_j \mathbf{e}_j \right) = \sum_{j=1}^n v_j f'(\mathbf{x}_0) \mathbf{e}_j = \sum_{j=1}^n v_j \frac{\partial f}{\partial x_j}(\mathbf{x}_0) $$
    This can be naturally represented in matrix multiplication form (where the derivative is a $m \times n$ matrix):
    $$ \begin{pmatrix}
    | & | & & | \\
    \frac{\partial f}{\partial x_1}(\mathbf{x}_0) & \frac{\partial f}{\partial x_2}(\mathbf{x}_0) & \cdots & \frac{\partial f}{\partial x_n}(\mathbf{x}_0) \\
    | & | & & |
    \end{pmatrix}_{m \times n}
    \begin{pmatrix}
    v_1 \\
    v_2 \\
    \vdots \\
    v_n
    \end{pmatrix}_{n \times 1}. $$
\end{enumerate}



\begin{definition}[Frobenius Norm]
Let $A=(a_{ij})_{i,j}\in M_{m\times n}(\mathbb{R})$. We can regard $A$ as a vector in $\mathbb{R}^{mn}$, and define its
Frobenius norm by
$$\|A\|_F = \sqrt{\sum_{i=1}^m \sum_{j=1}^n a_{ij}^2}.$$
\end{definition}

\begin{remark}
The Frobenius norm can also be elegantly expressed using the trace of a matrix. 
For $A \in M_{m \times n}(\mathbb{R})$, consider the matrix product $A A^\top$. By the definition of matrix multiplication, the $i$-th diagonal entry of $A A^\top$ is given by:
$$(A A^\top)_{ii} = \sum_{j=1}^n a_{ij} (A^\top)_{ji} = \sum_{j=1}^n a_{ij}^2$$
Since the trace of a matrix is the sum of its diagonal entries, taking the trace of $A A^\top$ yields:
$$\text{tr}(A A^\top) = \sum_{i=1}^m (A A^\top)_{ii} = \sum_{i=1}^m \sum_{j=1}^n a_{ij}^2$$
Therefore, the Frobenius norm can be equivalently written as:
$$\|A\|_F = \sqrt{\text{tr}(A A^\top)}.$$
\end{remark}


\begin{property}
Let $A, B \in M_{n \times n}(\mathbb{R})$ and $c \in \mathbb{R}$. The Frobenius norm is defined as $\|A\|_F = \sqrt{\sum_{i,j} a_{ij}^2} = \sqrt{\operatorname{tr}(A A^T)}$. It satisfies the following:
\begin{enumerate}
    \item \textbf{Positive Definiteness}: $\|A\|_F \ge 0$, and $\|A\|_F = 0 \iff A = \mathbf{0}$.
    \item \textbf{Absolute Homogeneity}: $\|c A\|_F = |c| \|A\|_F$.
    \item \textbf{Triangle Inequality}: $\|A + B\|_F \le \|A\|_F + \|B\|_F$.
    \item \textbf{Sub-multiplicativity}: $\|AB\|_F \le \|A\|_F \|B\|_F$. 
    \item \textbf{Power Inequality}: $\|A^k\|_F \le \|A\|_F^k$ for $k \in \mathbb{N}$.
    \item \textbf{Continuity}: For any column vector $\mathbf{x} \in \mathbb{R}^n$, $\|A\mathbf{x}\| \le \|A\|_F \|\mathbf{x}\|$. This implies that the linear map $\mathbf{x} \mapsto A\mathbf{x}$ is continuous on $\mathbb{R}^n$.
\end{enumerate}
\end{property}

\begin{pfbox}
We will only prove the sub-multiplicativity, and the rest is left as an exercise.
By the definition of matrix multiplication, the $(i, j)$-th entry of the product matrix $AB$ is given by:
$$(AB)_{ij} = \sum_{k=1}^n a_{ik} b_{kj}$$

Applying the Cauchy-Schwarz inequality to the vectors $(a_{i1}, \dots, a_{in})$ and $(b_{1j}, \dots, b_{nj})$, we obtain an upper bound for the square of this entry:
$$|(AB)_{ij}|^2 = \left| \sum_{k=1}^n a_{ik} b_{kj} \right|^2 \le \left( \sum_{k=1}^n a_{ik}^2 \right) \left( \sum_{k=1}^n b_{kj}^2 \right)$$

Now, we compute the squared Frobenius norm of the matrix $AB$ by summing over all entries $(1 \le i \le m$ and $1 \le j \le p)$:
\begin{align*}
\|AB\|_F^2 &= \sum_{i=1}^m \sum_{j=1}^p (AB)_{ij}^2 \\
&\le \sum_{i=1}^m \sum_{j=1}^p \left( \sum_{k=1}^n a_{ik}^2 \right) \left( \sum_{k=1}^n b_{kj}^2 \right)
\end{align*}

Notice that the sum over $i$ only depends on the terms $a_{ik}$, and the sum over $j$ only depends on the terms $b_{kj}$. Therefore, we can separate and factor the summations:
\begin{align*}
\|AB\|_F^2 &\le \left( \sum_{i=1}^m \sum_{k=1}^n a_{ik}^2 \right) \left( \sum_{j=1}^p \sum_{k=1}^n b_{kj}^2 \right) \\
&= \|A\|_F^2 \cdot \|B\|_F^2
\end{align*}

Finally, taking the square root of both sides, we get:
$$ \|AB\|_F \le \|A\|_F \|B\|_F $$
This completes the proof.
\end{pfbox}

\begin{theorem}[Neumann Series]
Let $M \in \mathbb{R}^{n \times n}$. If $\|M\|_F < 1$, then $(I + M)$ is invertible, and its inverse is given by the convergent series:
\[ (I + M)^{-1} = \sum_{i=0}^{\infty} (-1)^i M^i \]
\end{theorem}

\begin{pfbox}
Let $S_k = \sum_{i=0}^{k} (-1)^i M^i$. Since $(\mathbb{R}^{n \times n}, \|\cdot\|_F)$ is a complete metric space, it suffices to show that $\{S_k\}$ is a Cauchy sequence. 

For $k > m$:
\[ \|S_k - S_m\|_F = \left\| \sum_{i=m+1}^{k} (-1)^i M^i \right\|_F \le \sum_{i=m+1}^{k} \|M\|_F^i \le \frac{\|M\|_F^{m+1}}{1 - \|M\|_F} \]
As $m \to \infty$, the right-hand side vanishes since $\|M\|_F < 1$. Thus, $\{S_k\}$ is Cauchy and converges to some matrix $K$.

To verify $K = (I+M)^{-1}$, observe:
\begin{align*}
(I+M) S_k &= (I+M) \sum_{i=0}^{k} (-1)^i M^i = \sum_{i=0}^{k} (-1)^i M^i + \sum_{i=0}^{k} (-1)^i M^{i+1}\\
&=\sum_{i=0}^{k} (-1)^i M^i + \sum_{j=1}^{k+1} (-1)^{j-1} M^{j} = \sum_{i=0}^{k} (-1)^i M^i - \sum_{j=1}^{k+1} (-1)^j M^{j}.
\end{align*} 

Canceling the telescoping terms, we obtain:
\[ (I+M) S_k = I - (-1)^{k+1} M^{k+1} \]

Taking the limit as $k \to \infty$ (since $\|M\|_F < 1 \implies M^{k+1} \to \mathbf{0}$):
\[ (I+M) K = I + \mathbf{0} = I \]
Therefore, $K = (I+M)^{-1}$.
\end{pfbox}




\begin{theorem}[Chain Rule]
Let $E \subseteq \mathbb{R}^n,\, F\subseteq\mathbb{R}^m,\, f: E \to F,\, g: F\to\mathbb{R}^p$ and $\mathbb{x}_0\in int(E).$ 
Suppose $f$ is differentiable at $\mathbf{x}_0$, $f(\mathbf{x}_0)$ lies in the interior of $E$, and $g$ is differentiable at $f(\mathbf{x}_0)$.
Then $g \circ f$ is differentiable at $\mathbf{x}_0$, and
$$(g \circ f)'(\mathbf{x}_0) = g'(f(\mathbf{x}_0)) f'(\mathbf{x}_0).$$
\end{theorem}

\begin{pfbox}

Let $\mathbf{y}_0 = f(\mathbf{x}_0)$ and denote the derivatives by $A = f'(\mathbf{x}_0): \mathbb{R}^n \to \mathbb{R}^m$ and $B = g'(f(\mathbf{x}_0)) = g'(\mathbf{y}_0): \mathbb{R}^m \to \mathbb{R}^p$.

\vspace{0.8em}
\noindent \textbf{Step 1: Define the error function $\eta$} \\
Define $\eta: \mathbb{R}^m \to \mathbb{R}^p$ by:
$$
\eta(\mathbf{k}) = 
\begin{cases}
\frac{g(\mathbf{y}_0+\mathbf{k}) - g(\mathbf{y}_0) - B\mathbf{k}}{\|\mathbf{k}\|} & \text{if } \mathbf{k} \ne \mathbf{0} \\
\mathbf{0} & \text{if } \mathbf{k} = \mathbf{0}
\end{cases}
$$ 

Because $g$ is differentiable at $\mathbf{y}_0$, $\lim_{\mathbf{k} \to \mathbf{0}} \eta(\mathbf{k}) = \mathbf{0}$, which implies $\eta$ is continuous at $\mathbf{0}$.
Moreover, the relation
$$g(\mathbf{y}_0+\mathbf{k}) - g(\mathbf{y}_0) - B\mathbf{k} = \eta(\mathbf{k}) \|\mathbf{k}\|$$
is true for all $\mathbf{k} \in \mathbb{R}^m$.

\vspace{0.8em}
\noindent \textbf{Step 2: Expand the composition difference} \\
Take $\mathbf{h}$ such that $\mathbf{x}_0+\mathbf{h} \in E$. Let $\mathbf{k} = f(\mathbf{x}_0+\mathbf{h}) - f(\mathbf{x}_0)$.
Consider the difference:
\begin{align*}
(g \circ f)(\mathbf{x}_0+\mathbf{h}) - (g \circ f)(\mathbf{x}_0) 
&= g(f(\mathbf{x}_0+\mathbf{h})) - g(f(\mathbf{x}_0)) \\
&= g(\mathbf{y}_0+\mathbf{k}) - g(\mathbf{y}_0) \\
&= B(f(\mathbf{x}_0+\mathbf{h}) - f(\mathbf{x}_0)) \\
&\quad + \eta(f(\mathbf{x}_0+\mathbf{h}) - f(\mathbf{x}_0)) \|f(\mathbf{x}_0+\mathbf{h}) - f(\mathbf{x}_0)\|
\end{align*}

Subtracting $BA\mathbf{h}$ from both sides, we get:
\begin{align*}
&(g \circ f)(\mathbf{x}_0+\mathbf{h}) - (g \circ f)(\mathbf{x}_0) - BA\mathbf{h} \\
&= B(f(\mathbf{x}_0+\mathbf{h}) - f(\mathbf{x}_0) - A\mathbf{h}) + \eta(f(\mathbf{x}_0+\mathbf{h}) - f(\mathbf{x}_0)) \|f(\mathbf{x}_0+\mathbf{h}) - f(\mathbf{x}_0)\|
\end{align*} 

\vspace{0.8em}
\noindent \textbf{Step 3: Estimate the first term} \\
Note that:
$$\frac{\|B(f(\mathbf{x}_0+\mathbf{h}) - f(\mathbf{x}_0) - A\mathbf{h})\|}{\|\mathbf{h}\|} \le \|B\| \frac{\|f(\mathbf{x}_0+\mathbf{h}) - f(\mathbf{x}_0) - A\mathbf{h}\|}{\|\mathbf{h}\|}$$ 

By the assumption that $f$ is differentiable at $\mathbf{x}_0$, $\lim_{\mathbf{h} \to \mathbf{0}} \frac{\|f(\mathbf{x}_0+\mathbf{h}) - f(\mathbf{x}_0) - A\mathbf{h}\|}{\|\mathbf{h}\|} = 0$. By the Squeeze Theorem, the first term converges to zero.

\vspace{0.8em}
\noindent \textbf{Step 4: Estimate the second term} \\
Since $f$ is differentiable at $\mathbf{x}_0$, $f$ is continuous at $\mathbf{x}_0$, meaning $\lim_{\mathbf{h} \to \mathbf{0}} (f(\mathbf{x}_0+\mathbf{h}) - f(\mathbf{x}_0)) = \mathbf{0}$.
On the other hand, $\eta$ is continuous at $\mathbf{0}$ (and $\lim_{\mathbf{k} \to \mathbf{0}} \eta(\mathbf{k}) = \mathbf{0}$), it follows that $\lim_{\mathbf{h} \to \mathbf{0}} \eta(f(\mathbf{x}_0+\mathbf{h}) - f(\mathbf{x}_0)) = \mathbf{0}$

Next, consider:
\begin{align*}
\frac{\|f(\mathbf{x}_0+\mathbf{h}) - f(\mathbf{x}_0)\|}{\|\mathbf{h}\|} &= \frac{\|f(\mathbf{x}_0+\mathbf{h}) - f(\mathbf{x}_0) - A\mathbf{h} + A\mathbf{h}\|}{\|\mathbf{h}\|} \\
&\le \frac{\|f(\mathbf{x}_0+\mathbf{h}) - f(\mathbf{x}_0) - A\mathbf{h}\|}{\|\mathbf{h}\|} + \frac{\|A\mathbf{h}\|}{\|\mathbf{h}\|} \\
&\le \underbrace{\frac{\|f(\mathbf{x}_0+\mathbf{h}) - f(\mathbf{x}_0) - A\mathbf{h}\|}{\|\mathbf{h}\|}}_{\to 0\text{ by differentiability}} + \underbrace{\|A\|}_{fixed}
\end{align*} 

By the Squeeze Theorem:
$$\lim_{\mathbf{h} \to \mathbf{0}} \frac{\|\eta(f(\mathbf{x}_0+\mathbf{h}) - f(\mathbf{x}_0))\| \|f(\mathbf{x}_0+\mathbf{h}) - f(\mathbf{x}_0)\|}{\|\mathbf{h}\|} = 0$$

Combining Step 3 and Step 4, we conclude:
$$\lim_{\mathbf{h} \to \mathbf{0}} \frac{\|(g \circ f)(\mathbf{x}_0+\mathbf{h}) - (g \circ f)(\mathbf{x}_0) - BA\mathbf{h}\|}{\|\mathbf{h}\|} = 0$$ 
This proves that $g \circ f$ is differentiable at $\mathbf{x}_0$ with derivative $BA = g'(f(\mathbf{x}_0)) f'(\mathbf{x}_0)$.

\end{pfbox}

\begin{definition}
    Let $E\subseteq\mathbf{R}^n$ be open and $f:E\to\mathbf{R}^m.$
    We say $f$ is twice contiuously differentiable on $E$, denoted by $f\in C^2(E)$.
    \begin{enumerate}
    \item $f$ is continuously differentiable (its first order partial derivatives exist and are continuous).
    \item For each $i=1,2,\cdots n,$ the partial derivative $\frac{\partial f}{\partial x_i}: E\to\mathbf{R}^m$ is continously differentiable on $E$ (its second order partial derivatives exist and are continous on $E$).
\end{enumerate}
\end{definition}


\begin{theorem}[Clairaut]
Let $E$ be an open set and $f \in C^2(E)$. Then for all $1 \le i, j \le n$ and all $\mathbf{x}_0 \in E$:
$$\frac{\partial^2 f}{\partial x_i \partial x_j}(\mathbf{x}_0) = \frac{\partial^2 f}{\partial x_j \partial x_i}(\mathbf{x}_0).$$
\end{theorem}

\begin{pfbox}

Let $f = (f_1, f_2, \dots, f_m)$ where $f_k: E \to \mathbb{R}$. Without loss of generality, we can assume $f: E \to \mathbb{R}$ and the point of interest is $\mathbf{x}_0 = \mathbf{0}$. (By some shifting, $F(\mathbf{x}) = f(\mathbf{x} + \mathbf{x}_0) \Rightarrow F(\mathbf{0}) = f(\mathbf{x}_0)$). We also assume $i \neq j$.

Since $f \in C^2(E)$ and $\mathbf{0} \in E$, given any $\epsilon > 0$, there exists $\delta_0 > 0$ such that:
$$ \left| \frac{\partial^2 f}{\partial x_j \partial x_i}(\mathbf{x}) - \frac{\partial^2 f}{\partial x_j \partial x_i}(\mathbf{0}) \right| \le \frac{\epsilon}{2} \quad \text{and} \quad \left| \frac{\partial^2 f}{\partial x_i \partial x_j}(\mathbf{x}) - \frac{\partial^2 f}{\partial x_i \partial x_j}(\mathbf{0}) \right| \le \frac{\epsilon}{2} $$
if $\|\mathbf{x}\| < 2\delta_0$.

\vspace{0.8em}
\noindent \textbf{Step 1: The Second-Order Difference (DiD Approximation)}\\
Fix $\delta < \delta_0$, and consider the rectangle. Define the second-order difference $X$:
$$ X = f(\delta \mathbf{e}_i + \delta \mathbf{e}_j) - f(\delta \mathbf{e}_i) - f(\delta \mathbf{e}_j) + f(\mathbf{0}) $$
We can group the terms in two different ways:
\begin{align}
X &= \left( f(\delta \mathbf{e}_i + \delta \mathbf{e}_j) - f(\delta \mathbf{e}_i) \right) - \left( f(\delta \mathbf{e}_j) - f(\mathbf{0}) \right) \tag{1} \\
X &= \left( f(\delta \mathbf{e}_i + \delta \mathbf{e}_j) - f(\delta \mathbf{e}_j) \right) - \left( f(\delta \mathbf{e}_i) - f(\mathbf{0}) \right) \tag{2}
\end{align}

\vspace{0.8em}
\noindent \textbf{Step 2: Apply Mean Value Theorem for Grouping (1)}\\
Let $\phi(t) = f(t \mathbf{e}_i + \delta \mathbf{e}_j) - f(t \mathbf{e}_i)$. Then from (1), we have:
$$ X = \phi(\delta) - \phi(0) $$
By the MVT, $X = \delta \phi'(c_i)$ for some $c_i \in (0, \delta)$ \quad (*).

Note that $\phi'(t) = \frac{\partial f}{\partial x_i}(t \mathbf{e}_i + \delta \mathbf{e}_j) - \frac{\partial f}{\partial x_i}(t \mathbf{e}_i)$.
So, $\phi'(c_i) = \frac{\partial f}{\partial x_i}(c_i \mathbf{e}_i + \delta \mathbf{e}_j) - \frac{\partial f}{\partial x_i}(c_i \mathbf{e}_i)$.

Now, let $g(s) = \frac{\partial f}{\partial x_i}(c_i \mathbf{e}_i + s \mathbf{e}_j)$. Then $g'(s) = \frac{\partial^2 f}{\partial x_j \partial x_i}(c_i \mathbf{e}_i + s \mathbf{e}_j)$.
We can rewrite (*) as:
$$ \phi'(c_i) = g(\delta) - g(0) $$
By applying MVT again to $g(s)$, there exists $c_j \in (0, \delta)$ such that $g(\delta) - g(0) = \delta g'(c_j)$.
Thus, $\phi'(c_i) = \delta \frac{\partial^2 f}{\partial x_j \partial x_i}(c_i \mathbf{e}_i + c_j \mathbf{e}_j)$ \quad (**).

From (*) and (**), we obtain:
$$ X = \delta^2 \frac{\partial^2 f}{\partial x_j \partial x_i}(c_i \mathbf{e}_i + c_j \mathbf{e}_j) \quad \text{where } 0 < c_i < \delta, \ 0 < c_j < \delta. $$

\vspace{0.8em}
\noindent \textbf{Step 3: Symmetric Argument for Grouping (2)}\\
Similarly, starting from (2) and taking derivatives first with respect to $x_j$ and then $x_i$, MVT gives us:
$$ X = \delta^2 \frac{\partial^2 f}{\partial x_i \partial x_j}(d_i \mathbf{e}_i + d_j \mathbf{e}_j) \quad \text{where } 0 < d_i < \delta, \ 0 < d_j < \delta. $$

\vspace{0.8em}
\noindent \textbf{Step 4: Bounding the Error using Continuity}\\
Notice that the norm of the points we found is strictly bounded:
$$ \|c_i \mathbf{e}_i + c_j \mathbf{e}_j\| = \sqrt{c_i^2 + c_j^2} < \sqrt{\delta^2 + \delta^2} = \sqrt{2}\delta < \sqrt{2}\delta_0 < 2\delta_0 $$
Because these points fall within the $2\delta_0$ radius, our initial continuity conditions apply:
$$ \left| \frac{\partial^2 f}{\partial x_j \partial x_i}(c_i \mathbf{e}_i + c_j \mathbf{e}_j) - \frac{\partial^2 f}{\partial x_j \partial x_i}(\mathbf{0}) \right| \le \frac{\epsilon}{2} \implies \left| \frac{X}{\delta^2} - \frac{\partial^2 f}{\partial x_j \partial x_i}(\mathbf{0}) \right| \le  \frac{\epsilon}{2} $$
Similarly,
$$ \left| \frac{X}{\delta^2} - \frac{\partial^2 f}{\partial x_i \partial x_j}(\mathbf{0}) \right| \le  \frac{\epsilon}{2} $$

By the triangle inequality:
$$ \left| \frac{\partial^2 f}{\partial x_i \partial x_j}(\mathbf{0}) - \frac{\partial^2 f}{\partial x_j \partial x_i}(\mathbf{0}) \right| \le \epsilon \quad \forall \epsilon > 0 $$
Hence, we conclude that:
$$ \frac{\partial^2 f}{\partial x_i \partial x_j}(\mathbf{0}) = \frac{\partial^2 f}{\partial x_j \partial x_i}(\mathbf{0}). $$

\end{pfbox}

\begin{rmk}[Numerical Approximation via Second-Order Difference]
From the proof above, we know that by the continuity of $\frac{\partial^2 f}{\partial x_j \partial x_i}$ and $\frac{\partial^2 f}{\partial x_i \partial x_j}$,
$$ \lim_{\delta \to 0} \frac{X}{\delta^2} = \frac{\partial^2 f}{\partial x_j \partial x_i}(\mathbf{0}) \quad \text{and} \quad \lim_{\delta \to 0} \frac{X}{\delta^2} = \frac{\partial^2 f}{\partial x_i \partial x_j}(\mathbf{0}) $$
Hence, if $f \in C^2(E)$, for any point $\mathbf{x}_0$, we can use numerical approximation by a second-order finite difference without computing the derivative directly:
$$ \lim_{\delta \to 0} \frac{f(\mathbf{x}_0 + \delta \mathbf{e}_i + \delta \mathbf{e}_j) - f(\mathbf{x}_0 + \delta \mathbf{e}_i) - f(\mathbf{x}_0 + \delta \mathbf{e}_j) + f(\mathbf{x}_0)}{\delta^2} = \frac{\partial^2 f}{\partial x_j \partial x_i}(\mathbf{x}_0) = \frac{\partial^2 f}{\partial x_i \partial x_j}(\mathbf{x}_0) $$
\end{rmk}


\subsection{Inverse Function Theorem}

\begin{definition}[Metric Space]
The pair $(X,d)$ is called a metric space if a set $X$ endowed with the function $d:X\times X\to\mathbb{R}$ such that $\dots$
\end{definition}

\begin{definition}[Convergence]
Given a sequence $\{x_n\}_{n=1}^\infty$ in $X$. We say the sequence converges to $x$ if $\forall \varepsilon > 0\,\exists N \in \mathbb{N}$ s.t. for $n \ge N$, $d(x_n, x) < \varepsilon$.
\end{definition}

\begin{definition}[Cauchy Sequence]
A sequence $\{x_n\}_{n=1}^\infty$ is said to be Cauchy if $\forall \varepsilon > 0$, $\exists N \in \mathbb{N}$ s.t. for $m,n \ge N$, $d(x_n, x_m) < \varepsilon$.
\end{definition}

\begin{definition}[Completeness]
$(X,d)$ is complete if every Cauchy sequence in $X$ converges to a limit in $X$.
\end{definition}

\begin{theorem}
Let $E$ be a closed subset in $X$, and $(X,d)$ is a complete metric space. Then $(E,d)$ is also complete.
\end{theorem}

\begin{pfbox}

    Take any Cauchy sequence in $E$, which is also a Cauchy in $X$. By completeness of $X$, the sequence converges to some point in $X$. Since a closed set contains all limit point, the Cauchy sequence converges to a limit in $E$.

\end{pfbox}

\begin{definition}[Contraction]
Let $(X,d)$ be a metric space and $T: X \to X$ be a map.
\begin{enumerate}
    \item We say $T$ is a contraction if $d(T(x), T(y)) \le d(x,y)$ $\forall x, y \in X$.
    \item $T$ is a strict contraction if there exists a constant $0<c<1$ s.t. $d(T(x), T(y)) \le cd(x,y)$.
\end{enumerate}
\end{definition}

\begin{remark}
Under the standard setting, the term \textbf{contraction} simply refers to a \textbf{strict contraction}. Terence Tao use ``contraction'' for maps where $d(T(x), T(y)) \le d(x,y)$ and explicitly write ``strict contraction'' when the constant is strictly less than $1$.
\end{remark}

\begin{definition}[Fixed Point]
Let $T: X \to X$ be a map and $x^* \in X$. We say $x^*$ is a fixed point of $T$ if $T(x^*) = x^*$.
\end{definition}

\begin{remark}
Suppose $X$ is a subset of $\mathbb{R}$. The existence of a fixed point means that the graph of $y=f(x)$ intersect with $y=x$.
\end{remark}


\begin{theorem}[Contraction Mapping Theorem]
Let $(X,d)$ be a metric space and $T: X \to X$ be a (strict) contraction. Then $T$ has at most one fixed point. Moreover, if we assume $X$ is non-empty and complete, then $T$ has exactly one fixed point.
\end{theorem}

\begin{pfbox}

Suppose $x,y$ are distinct fixed point of $T$,
consider $d(x,y)=d(T(x),T(y))\le cd(x,y)$ for some $c\in (0,1).$ 
The inequality $(1-c)d(x,y)\le 0$ forces $x=y.$\\
Now we further assume $X$ is complete. Take any point $x_0$, construct the sequence $\{x_n\}^\infty_{n=0},\, x_{n+1}=T(x_n)\,\forall n\in\mathbb{N}.$
Observe that $d(x_{n+1}, x_n)\le c^n d(x_1,x_0)$ for $n\ge 1$.
Suppose $n > m \ge 1$,
\begin{align*}
    d(x_n, x_m) &\le d(x_n, x_{n-1}) + d(x_{n-1}, x_{n-2}) + \cdots + d(x_{m+1}, x_m)\\
    &\le c^{n-1} d(x_1, x_0) + c^{n-2} d(x_1, x_0) + \cdots + c^m d(x_1, x_0) \\
    &\le \frac{c^m}{1-c}d(x_1, x_0).
\end{align*}
Take $m \to \infty$, $c^m \to 0$, which shows $\{x_n\}_{n=1}^\infty$ is a Cauchy in $X$. Since $X$ is complete, $\lim_{n \to \infty} x_n$ exists and converges to $x^*$.
Finally,
\begin{align*}
    d(x^*, T(x^*)) &\le d(x^*, x_{n+1}) + d(x_{n+1}, T(x^*)) \\
    &\le  d(x^*, x_{n+1}) + d(T(x_n), T(x^*))\\
    &\le d(x^*, x_{n+1}) + cd(x_n, x^*). 
\end{align*}
This shows $x^*=T(x^*).$

\end{pfbox}


\begin{definition}[Lipschitz continuous]
    Given $(X,d_X), (Y,d_Y)$ and a map $f:X\to Y.$
    We say $f$ is Lipschitz continuous if there exists a constant $L\ge 0$ such that
    $d_Y(f(x_1),f(x_2))\le Ld_X(x_1,x_2)$ for any $x_1,x_2\in X.$
    Any such $L$ is referred to as a Lipschitz constant for the function $f$, 
    and $f$ may also be referred to as $L-$Lipschitz.
\end{definition}

\begin{definition}[Uniform Continuity]
    Let $f: X\rightarrow Y$ be a map btw $(X,d_X)$ and $(Y,d_Y)$.
    $f$ is uniformly continous if $\forall \epsilon>0\,\exists\,\delta>0$ s.t. for any $x,y\in X$ $d_X(x,y)<\delta\rightarrow d_Y(f(x),f(y))<\epsilon.$
    This means $\delta$ is independent of the location of $x\in X$.
\end{definition}

\begin{theorem}
    If a function $f$ is Lipschitz continuous, then
    \begin{enumerate}
        \item $f$ is continous.
        \item $f$ is uniformly continous.
    \end{enumerate}
\label{lip_conti implies conti and uniform conti}
\end{theorem}

\begin{pfbox}
    Let $f:X\to Y$ be a $L-$Lipschitz continous map. Fix $\epsilon>0$ and a point $a\in X,$ choose $\frac{\varepsilon}{L}$ such that
    if $d(x,a)<\frac{\varepsilon}{L},$ one has $d(f(x),f(a))\le Ld(x,a)<L\frac{\varepsilon}{L}=\varepsilon,$ which shows that $f$ is contiuous at $a$. 
    The argument for 2 is analogous.
\end{pfbox}

\begin{corollary}
A contraction $T:X\to X$ is a continuous function.
\end{corollary}

\begin{pfbox}
    A contraction with contraction ratio $c$ is $c-$Lipschitz, by theorem~\eqref{lip_conti implies conti and uniform conti}, $T$ is continous.
\end{pfbox}

\begin{theorem}
    Let $T: \mathbb{R}^n \to \mathbb{R}^m$ be a linear transformation. Then $T$ is continuous.
\label{linear_map_continuous}
\end{theorem}

\begin{pfbox}
    Let $\{e_1, e_2, \dots, e_n\}$ be the standard basis for $\mathbb{R}^n$. Any vector $x \in \mathbb{R}^n$ can be expressed as $x = \sum_{i=1}^n x_i e_i$. By the linearity of $T$ and the triangle inequality, we have:
    $$ \|T(x)\| = \left\|T\left(\sum_{i=1}^n x_i e_i\right)\right\| = \left\|\sum_{i=1}^n x_i T(e_i)\right\| \le \sum_{i=1}^n |x_i| \|T(e_i)\| $$
    
    In the standard Euclidean norm, the absolute value of any component is bounded by the norm of the vector, so $|x_i| \le \|x\|$ for all $i=1, \dots, n$. Substituting this yields:
    $$ \|T(x)\| \le \sum_{i=1}^n \|x\| \|T(e_i)\| = \|x\| \sum_{i=1}^n \|T(e_i)\| $$
    
    Define $M = \sum_{i=1}^n \|T(e_i)\|$. Since $\{T(e_i)\}$ is a finite set of fixed vectors in $\mathbb{R}^m$, $M$ is a finite non-negative constant independent of $x$. Thus, we obtain the boundedness condition $\|T(x)\| \le M\|x\|$. 
    
    For any $x, y \in \mathbb{R}^n$, the linearity of $T$ implies:
    $$ \|T(x) - T(y)\| = \|T(x - y)\| \le M\|x - y\| $$
    
    This demonstrates that $T$ is $M$-Lipschitz continuous. By theorem~\eqref{lip_conti implies conti and uniform conti}, it follows that $T$ is continuous.
\end{pfbox}

\begin{lemma}
Let $B(\mathbf{0},r)$ be an open ball in $\mathbb{R}^n$ and $g: B(\mathbf{0},r) \to \mathbb{R}^n$ be a map with $g(\mathbf{0}) = \mathbf{0}$. 
If $\|g(\mathbf{x}) - g(\mathbf{y})\| \le \frac{1}{2}\|\mathbf{x} - \mathbf{y}\|$ for $\mathbf{x}, \mathbf{y} \in B(\mathbf{0},r)$, then the function $f: B(\mathbf{0},r) \to \mathbb{R}^n$ defined by $f(\mathbf{x}) = \mathbf{x} + g(\mathbf{x})$ is one-to-one, and the image of $f(B(\mathbf{0},r))$ contains $B(\mathbf{0}, \frac{r}{2})$.
\label{lemma_for_inverse_funct_thm}
\end{lemma}

\begin{pfbox}

We first show injection. Suppose $f(\mathbf{x}) = f(\mathbf{y}),$ by construction, we have $\|\mathbf{x} - \mathbf{y}\| = \|g(\mathbf{x}) - g(\mathbf{y})\| \le \frac{1}{2}\|\mathbf{x} - \mathbf{y}\|.$
The inequality $\frac{1}{2}\|\mathbf{x} - \mathbf{y}\| \le 0$ forces $\|\mathbf{x} - \mathbf{y}\| = 0$. So $f$ is one-to-one.\\
Next, we want to show $B(\mathbf{0}, \frac{r}{2}) \subseteq f(B(\mathbf{0},r))$. It suffices to show $\forall\,\mathbf{y} \in B(\mathbf{0}, \frac{r}{2})$, $\exists\,\mathbf{x} \in B(\mathbf{0},r)$ s.t.
$f(\mathbf{x}) = \mathbf{y}.$ Define $F(\mathbf{x}) = \mathbf{y} - g(\mathbf{x})$. If we can find the fixed point of $F$, then we equivalently find $\mathbf{y} = f(\mathbf{x})$.
Fix $\mathbf{y} \in B(\mathbf{0}, \frac{r}{2})$. $\|\mathbf{y}\| < \frac{r}{2}$. $\exists\,\varepsilon > 0$ s.t. $\|\mathbf{y}\| \le \frac{r-\varepsilon}{2}$.
Consider the map $F: \overline{B(\mathbf{0}, r-\varepsilon)} \to \overline{B(\mathbf{0}, r-\varepsilon)}.$ Take $\mathbf{x} \in \overline{B(\mathbf{0}, r-\varepsilon)}, \,\|\mathbf{x}\| \le r - \varepsilon$. 
Then we have $\|F(\mathbf{x})\| \le \|\mathbf{y}\| + \|g(\mathbf{x})\| \le \frac{r-\varepsilon}{2} + \|g(\mathbf{x})\| \le \frac{r-\varepsilon}{2} + \frac{r-\varepsilon}{2} = r-\varepsilon,$
where $\|g(\mathbf{x})\| = \|g(\mathbf{x}) - g(\mathbf{0})\| \le \frac{1}{2}\|\mathbf{x}\| \le \frac{1}{2}(r-\varepsilon)$. So $F(\mathbf{x}) \in \overline{B(\mathbf{0}, r-\varepsilon)}$.
Given a fixed $\mathbf{y}$,
$\|F(\mathbf{x}_1) - F(\mathbf{x}_2)\| = \|\mathbf{y} - g(\mathbf{x}_1) - (\mathbf{y} - g(\mathbf{x}_2))\| \le \|g(\mathbf{x}_2) - g(\mathbf{x}_1)\| \le \frac{1}{2}\|\mathbf{x}_2 - \mathbf{x}_1\|.$ Hence, $F: \overline{B(\mathbf{0}, r-\varepsilon)} \to \overline{B(\mathbf{0}, r-\varepsilon)}$ is a contraction $\forall\,\mathbf{y} \in B(0, \frac{r}{2})$. 
Since $\overline{B(\mathbf{0}, r-\varepsilon)}$ is complete, $\exists\,\mathbf{x}^* \in \overline{B(\mathbf{0}, r-\varepsilon)}$ s.t. $\mathbf{x}^* = F(\mathbf{x}^*)$
This shows $F(\mathbf{x}^*) = \mathbf{y} - g(\mathbf{x}^*) = \mathbf{x}^*$, and hence $\mathbf{x}^* + g(\mathbf{x}^*) = \mathbf{y}.$

\end{pfbox}



\begin{theorem}[Inverse Function Theorem]
Let $E$ be open in $\mathbb{R}^n$ and $f:E\to\mathbb{R}^n\in C^1(E).$ Suppose the linear transformation $f'(\mathbf{x}_0):\mathbb{R}^n\to\mathbb{R}^n$ is invertible, then there exists an open set $U\subseteq E$ containing $\mathbf{x}_0$ and an open set $V\subseteq\mathbb{R}^n$ containing $f(\mathbf{x}_0)$ such that $f:U\to V$ is a bijection.
Moreover, the inverse map $f^{-1}:V\to U$ is also continuously differentiable, and 
$$(f^{-1})'(f(\mathbf{x_0}))=(f'(\mathbf{x}_0))^{-1}.$$
\end{theorem}


\begin{pfbox}

Let $A=f'(\mathbf{x}_0)$. Firstly, $f$ is differentiable at $\mathbf{x}_0$, so:
$$ \lim_{\mathbf{x}\to \mathbf{0}}\frac{\|f(\mathbf{x}_0+\mathbf{x})-f(\mathbf{x}_0)-A\mathbf{x}\|}{\|\mathbf{x}\|} = 0. $$
Given $A$ is invertible, as $\mathbf{x}\to\mathbf{0}$, we have:
$$ \frac{\|A^{-1}\big(f(\mathbf{x}_0+\mathbf{x})-f(\mathbf{x}_0)-A\mathbf{x}\big)\|}{\|\mathbf{x}\|} \le \|A^{-1}\|\frac{\|f(\mathbf{x}_0+\mathbf{x})-f(\mathbf{x}_0)-A\mathbf{x}\|}{\|\mathbf{x}\|} \to 0. $$

\vspace{0.8em}
\noindent \textbf{Step 1: Construct the auxiliary function $g$ and show it is a contraction}\\
Let $g(\mathbf{x})=A^{-1}\big(f(\mathbf{x}_0+\mathbf{x})-f(\mathbf{x}_0)\big)-\mathbf{x}$. It is clear that $g(\mathbf{0})=\mathbf{0}$.
Substituting this into our limit condition yields $\lim_{\mathbf{x}\to\mathbf{0}}\frac{\|g(\mathbf{x})\|}{\|\mathbf{x}\|}=0$. 
Rewrite the limit condition as:
$$ \lim_{\mathbf{x}\to\mathbf{0}}\frac{\|g(\mathbf{x})-g(\mathbf{0})-g'(\mathbf{0})\mathbf{x}\|}{\|\mathbf{x}\|} = 0, $$
which shows that $g$ is differentiable at $\mathbf{0}$ with $g'(\mathbf{0})=\mathbf{0}$.

Since both $f$ and $A$ are $C^1$, $g$ is also $C^1$. By definition, there exists an $r>0$ such that for all $\mathbf{x}\in B(\mathbf{0},r)$:
$$ \left\|\frac{\partial g}{\partial x_j}(\mathbf{x})-\frac{\partial g}{\partial x_j}(\mathbf{0})\right\| = \left\|\frac{\partial g}{\partial x_j}(\mathbf{x})\right\| \le \frac{1}{2n^2}. $$
For $\mathbf{x}\in B(\mathbf{0},r)$ and $\mathbf{v}\in \mathbb{R}^n$,
$$ \|D_{\mathbf{v}}g(\mathbf{x})\| = \left\|\sum_{j=1}^{n}v_j\frac{\partial g}{\partial x_j}(\mathbf{x})\right\| \le \sum_{j=1}^{n}|v_j|\left\|\frac{\partial g}{\partial x_j}(\mathbf{x})\right\| \le \sum_{j=1}^{n}\|\mathbf{v}\|\frac{1}{2n^2} = \frac{1}{2n}\|\mathbf{v}\| \le \frac{1}{2}\|\mathbf{v}\|. $$
For $\mathbf{x},\mathbf{y}\in B(\mathbf{0},r)$, consider:
\begin{align*}
g(\mathbf{y})-g(\mathbf{x}) &= \int_{0}^{1}\frac{d}{dt}g(\mathbf{x}+t(\mathbf{y}-\mathbf{x}))dt \quad \text{(FTC(1))} \\
&= \int_{0}^{1}Dg(\mathbf{x}+t(\mathbf{y}-\mathbf{x}))(\mathbf{y}-\mathbf{x})dt \quad \text{(Chain rule)} \\
&= \int_{0}^{1}D_{\mathbf{y}-\mathbf{x}}g(\mathbf{x}+t(\mathbf{y}-\mathbf{x}))dt \quad \text{(Since $g$ is $C^1$)}
\end{align*}
Therefore,
$$ \|g(\mathbf{y})-g(\mathbf{x})\| = \left\|\int_{0}^{1}D_{\mathbf{y}-\mathbf{x}}g(\mathbf{x}+t(\mathbf{y}-\mathbf{x}))dt\right\| \le \int_{0}^{1}\frac{1}{2}\|\mathbf{y}-\mathbf{x}\|dt \le \frac{1}{2}\|\mathbf{y}-\mathbf{x}\|, $$
which means $g$ is a contraction.

\vspace{0.8em}
\noindent \textbf{Step 2: Define $\Phi$, establish local bijection, and show $\Phi^{-1}$ is continuous}\\
Define $\Phi(\mathbf{x}) = g(\mathbf{x})+\mathbf{x} = A^{-1}\big(f(\mathbf{x}_0+\mathbf{x})-f(\mathbf{x}_0)\big)$, where $g(\mathbf{0})=\mathbf{0}$ and $g$ is a contraction.
Apply the lemma~\eqref{lemma_for_inverse_funct_thm}, $\Phi$ is one-to-one on $B(\mathbf{0},r)$ and $B(\mathbf{0},\frac{r}{2})\subseteq\Phi(B(\mathbf{0},r))$.
Set $W=\Phi^{-1}(B(\mathbf{0},\frac{r}{2}))$, then $\Phi|_W:W\to B(\mathbf{0},\frac{r}{2})$ is a bijection.

Clearly, $\Phi$ is continuous. We claim that $\Phi^{-1}$ is also continuous. For $\mathbf{x},\mathbf{y}\in B(\mathbf{0},r)$:
$$ \|\Phi(\mathbf{x})-\Phi(\mathbf{y})\| = \|(\mathbf{x}-\mathbf{y})+(g(\mathbf{x})-g(\mathbf{y}))\| \ge \|\mathbf{x}-\mathbf{y}\|-\|g(\mathbf{x})-g(\mathbf{y})\| \ge \frac{1}{2}\|\mathbf{x}-\mathbf{y}\|. $$
Let $\mathbf{x}=\Phi^{-1}(\mathbf{u})$ and $\mathbf{y}=\Phi^{-1}(\mathbf{v})$, with $\mathbf{u},\mathbf{v}\in B(\mathbf{0},\frac{r}{2})\subseteq \Phi(B(\mathbf{0},r))$, then:
\begin{align*}
\|\Phi(\Phi^{-1}(\mathbf{u}))-\Phi(\Phi^{-1}(\mathbf{v}))\| &\ge \frac{1}{2}\|\Phi^{-1}(\mathbf{u})-\Phi^{-1}(\mathbf{v})\| \\
\|\Phi^{-1}(\mathbf{u})-\Phi^{-1}(\mathbf{v})\| &\le 2\|\mathbf{u}-\mathbf{v}\|.
\end{align*}
This shows that $\Phi^{-1}$ is Lipschitz continuous and hence continuous.

\vspace{0.8em}
\noindent \textbf{Step 3: Prove $\Phi^{-1}$ is differentiable at $\mathbf{0}$}\\
We claim that the derivative of $\Phi^{-1}$ at $\mathbf{0}$ is the identity map $\text{id}$.
Let $\mathbf{z}\in B(\mathbf{0},\frac{r}{2})\setminus\{\mathbf{0}\}$. Since $\Phi|_W$ is bijective, let $\mathbf{z}=\Phi(\mathbf{x})=\mathbf{x}+g(\mathbf{x})$, which means $\mathbf{x}=\Phi^{-1}(\mathbf{z})$ and $\mathbf{x}\neq\mathbf{0}$.
We want to show:
$$ \lim_{\mathbf{z}\to\mathbf{0}}\frac{\|\Phi^{-1}(\mathbf{z})-\Phi^{-1}(\mathbf{0})-\text{id}(\mathbf{z})\|}{\|\mathbf{z}\|} = 0. $$
Since $\Phi^{-1}(\mathbf{0})=\mathbf{0}$, the expression inside the limit simplifies to:
$$ \frac{\|\Phi^{-1}(\mathbf{z})-\mathbf{z}\|}{\|\mathbf{z}\|} = \frac{\|\mathbf{x}-(\mathbf{x}+g(\mathbf{x}))\|}{\|\mathbf{z}\|} = \frac{\|-g(\mathbf{x})\|}{\|\mathbf{z}\|} = \frac{\|g(\mathbf{x})\|}{\|\mathbf{x}\|}\cdot\frac{\|\mathbf{x}\|}{\|\mathbf{z}\|}. $$
From the Lipschitz continuity in Step 2, taking $\mathbf{v}=\mathbf{0}$, we have $\|\Phi^{-1}(\mathbf{z})\| \le 2\|\mathbf{z}\|$, which implies $\frac{\|\mathbf{x}\|}{\|\mathbf{z}\|} \le 2$.
Furthermore, as $\mathbf{z}\to\mathbf{0}$, sequential continuity dictates that $\mathbf{x}=\Phi^{-1}(\mathbf{z})\to\mathbf{0}$.
Using the Squeeze Theorem:
$$ \lim_{\mathbf{z}\to\mathbf{0}}\frac{\|\Phi^{-1}(\mathbf{z})-\mathbf{z}\|}{\|\mathbf{z}\|} = \lim_{\mathbf{x}\to\mathbf{0}}\left(\frac{\|g(\mathbf{x})\|}{\|\mathbf{x}\|}\cdot\frac{\|\mathbf{x}\|}{\|\mathbf{z}\|}\right) \le \lim_{\mathbf{x}\to\mathbf{0}}\left(\frac{\|g(\mathbf{x})\|}{\|\mathbf{x}\|}\cdot 2\right) = 0 \cdot 2 = 0. $$
This proves that $\Phi^{-1}$ is differentiable at $\mathbf{0}$ with $(\Phi^{-1})'(\mathbf{0})=\text{id}$.

\vspace{0.8em}
\noindent \textbf{Step 4: Prove $f^{-1}$ is differentiable at $f(\mathbf{x}_0)$}\\
Recall the definition $\Phi(\mathbf{x}) = A^{-1}\big(f(\mathbf{x}_0+\mathbf{x})-f(\mathbf{x}_0)\big)$. Multiplying by $A$ and rearranging yields:
$$ f(\mathbf{x}+\mathbf{x}_0) = A\Phi(\mathbf{x}) + f(\mathbf{x}_0). $$
Since $\Phi$ is a local bijection, $f$ is also one-to-one and onto near $\mathbf{x}_0$. 
Let $\mathbf{y} = f(\mathbf{x}+\mathbf{x}_0)$, then $\mathbf{y} = A\Phi(\mathbf{x}) + f(\mathbf{x}_0)$. We can express $\mathbf{x}$ as:
$$ \mathbf{x} = \Phi^{-1}\big(A^{-1}(\mathbf{y}-f(\mathbf{x}_0))\big). $$
Thus, the inverse function $f^{-1}$ is given by:
$$ f^{-1}(\mathbf{y}) = \mathbf{x} + \mathbf{x}_0 = \Phi^{-1}\big(A^{-1}(\mathbf{y}-f(\mathbf{x}_0))\big) + \mathbf{x}_0. $$
We apply the Chain Rule to differentiate $f^{-1}$ at $f(\mathbf{x}_0)$:
\begin{align*}
(f^{-1})'(f(\mathbf{x}_0)) &= (\Phi^{-1})'\big(A^{-1}(f(\mathbf{x}_0)-f(\mathbf{x}_0))\big) \circ (A^{-1})'(f(\mathbf{x}_0)-f(\mathbf{x}_0)) \\
&= (\Phi^{-1})'(\mathbf{0}) \circ A^{-1}(\mathbf{0}) \\
&= \text{id} \circ A^{-1} = A^{-1} = (f'(\mathbf{x}_0))^{-1}.
\end{align*}

\end{pfbox}

\subsection{Implicit Function Theorem}

\begin{definition}[Graph]
Given a function $f: U\subseteq X \to Y$, the graph of $f$ is defined as the following subset of Cartesian product space:
$$\mathcal{G}(f) = \{(x,f(x)):x\in X\}.$$
\end{definition}

\begin{example}
Let $f: U \subseteq \mathbb{R}^{n-1} \to \mathbb{R}$ be a function. Define the graph of $f$ as a $(n-1)-$dimension subset in $\mathbb{R}^n$: 
$$\{ (x_1, \ldots, x_{n-1}, f(x_1, \ldots, x_{n-1})) \in \mathbb{R}^n \mid (x_1, \ldots, x_{n-1}) \in U \}.$$
\end{example}


\begin{definition}[Level Set]
Let $U \subseteq \mathbb{R}^n$ be an open set and $f: U \to \mathbb{R}$ be a real-valued function. For a given constant $c \in \mathbb{R}$, the set of all points $\mathbf{x} \in U$ that satisfy the equation $f(\mathbf{x}) = c$ forms a subset of $U$. This subset is formally called the \textbf{level set} of $f$ at level $c$, denoted by:
$$L_c = \{ \mathbf{x} \in U \mid f(\mathbf{x}) = c \}.$$
\end{definition}


\begin{example}
\begin{enumerate}
    \item \textbf{Level Curve ($n=2$):} A 1D level set in a 2D space is defined by a 2-variable equation $f(x,y)=c$. We call this a \textbf{level curve}/\textbf{contour line}/\textbf{isoline} of $f$ at level $c$.
    \item \textbf{Level Surface ($n=3$):} A 2D level set in a 3D space is defined by a 3-variable equation $f(x,y,z)=c$. We call this a \textbf{level surface}/\textbf{isosurface} of $f$ at level $c$.
    \item \textbf{Level Hypersurface ($n>3$):} An $(n-1)$-dimensional level set in an $n$-dimensional space is defined by $f(x_1, \dots, x_n)=c$. We simply call this a \textbf{level hypersurface}.
\end{enumerate}
\end{example}

\begin{remark}
    Let $U\subseteq\mathbb{R}^3$ and a three-variable function $f:U\to\mathbb{R}$ defined by $w=f(x,y,z)$.
    \begin{enumerate}
        \item The graph of $f$ is defined as $\{(x,y,z,f(x,y,z)):(x,y,z)\in U\},$ which is a 3-dimensional subset in $\mathbb{R}^4.$
        \item Set $w=c$. Given this constant $c$, we have a level surface of $f$ at level $c$, which is a 2-dimensional subset $S=\{(x,y,z)\in U: f(x,y,z)=c\}$ in $\mathbb{R}^3.$ 
        \item Locally, $S$ can be expressed as the graph of a 2-variable function, $S'=\{(x,y,z(x,y)):(x,y)\in\mathbb{R}^2\}$, which is a 2-dimensional subset in $\mathbb{R}^3,$ where $S'\subseteq S.$ 
    \end{enumerate}
\end{remark}

\begin{remark}
    A related but distinct concept is the intersection of two surfaces. For example, suppose we have a double cone and a cutting plane, both of which are level surfaces defined by 3-variable equations. 
    Their intersection forms a conic section (a quadratic curve). Intuitively, a system of 3 variables restricted by 2 independent equations reduces the degrees of freedom to 1 ($3 - 2 = 1$), resulting in a 1-dimensional curve. 
    However, this is not the same concept as a "level curve." A standard level curve is defined by a single equation in a 2D space (e.g., $f(x,y)=c$), whereas this intersection is a 1D curve living in a 3D space, defined by a system of two equations.
\end{remark}


\begin{theorem}[Implicit Function Theorem]
Let $E$ be an open set in $\mathbb{R}^n$, and $f:E\to\mathbb{R}$ is $C^1$. Suppose $y\in E$ such that $f(y)=0$ and $f_{x_n}=\frac{\partial f}{\partial x_n}\neq 0,$ then:
\begin{enumerate}
    \item There exists an open subset $U\subseteq\mathbb{R}^{n-1}$ containing $(y_1,y_2,\cdots,y_{n-1})$, and a map $g:U\to\mathbb{R}$ such that $g(y_1,y_2,\cdots,y_{n-1})=y_n.$
    \item The level set $\{x\in V: f(x)=0\}$ is a graph of a function of $g$ over $U$, i.e. $\{x\in V: f(x)=0\}=\{(x_1,x_2,\cdots,x_{n-1},g(x_1,x_2,\cdots,x_{n-1})):(x_1,x_2,\cdots,x_{n-1})\in U\}.$
    \item For $j=1,2,\cdots,n-1$,
    $$\frac{\partial g}{\partial x_j}(y_1, y_2,\cdots,y_{n-1})=-\frac{\frac{\partial f}{\partial x_j}(y)}{\frac{\partial f}{\partial x_n}(y)}.$$ 
\end{enumerate}

\end{theorem}


\begin{pfbox}
\vspace{0.8em}
\noindent \textbf{Step 1: Define a new function $F$ and check its Jacobian}\\
Define $F: E \to \mathbb{R}^n$ by $(x_1, x_2, \ldots, x_n) \mapsto (x_1, x_2, \ldots, x_{n-1}, f(x_1, x_2, \ldots, x_n))$.

Since $f \in C^1(E)$, it follows that $F \in C^1(E)$. The Jacobian matrix of $F$ at $x$ is:
$$ DF(x) = \begin{pmatrix} 
- & \nabla x_1 & - \\ 
- & \nabla x_2 & - \\ 
  & \vdots     &   \\ 
- & \nabla f   & - 
\end{pmatrix} 
= \begin{pmatrix} 
1 & 0 & \cdots & 0 \\ 
0 & 1 & \cdots & 0 \\ 
\vdots & \vdots & \ddots & \vdots \\ 
\frac{\partial f}{\partial x_1}(x) & \frac{\partial f}{\partial x_2}(x) & \cdots & \frac{\partial f}{\partial x_n}(x) 
\end{pmatrix} $$
We are given that $\frac{\partial f}{\partial x_n}(y) \neq 0$. Therefore, the determinant evaluated at $y$ is:
$$ \det(DF(y)) = \frac{\partial f}{\partial x_n}(y) \neq 0, $$
which means $DF(y)$ is invertible.

\vspace{0.8em}
\noindent \textbf{Step 2: Apply the Inverse Function Theorem to construct local inverse}\\
By the Inverse Function Theorem, there exists an open set $V \subseteq E$ and $W \subseteq \mathbb{R}^n$ such that $y \in V$, $F(y) \in W$, and $F: V \to W$ is a $C^1$ diffeomorphism. 
Write the inverse in coordinates as $F^{-1}(u) = (h_1(u), h_2(u), \ldots, h_n(u))$ for $u = (u_1, u_2, \ldots, u_n) \in W$.
By construction, $F(F^{-1}(u)) = u$, which implies:
$$ F(h_1(u), h_2(u), \ldots, h_n(u)) = (u_1, u_2, \ldots, u_n). $$
Using the definition of $F$, the left-hand side expands to:
$$ (h_1(u), h_2(u), \ldots, h_{n-1}(u), f(h_1(u), \ldots, h_n(u))). $$
Equating the components element by element, we obtain:
$$ h_1(u) = u_1, \quad h_2(u) = u_2, \quad \ldots, \quad h_{n-1}(u) = u_{n-1}, $$
$$ f(h_1(u), h_2(u), \ldots, h_n(u)) = u_n \implies f(u_1, u_2, \ldots, u_{n-1}, h_n(u)) = u_n\quad (*).$$
Thus, the inverse function takes the form:
$$ F^{-1}(u_1, u_2, \ldots, u_n) = (u_1, u_2, \ldots, u_{n-1}, h_n(u_1, u_2, \ldots, u_n)). $$

\vspace{0.8em}
\noindent \textbf{Step 3: Restrict to $u_n = 0$ and define the implicit function $g$}\\
If we restrict the domain to $u_n = 0$, by $(*),$ we have:
$$
f(u_1, \ldots, u_{n-1}, h_n(u_1, \ldots, u_{n-1}, 0)) = 0.
$$

Set an open set $U = \{ x' = (x_1, x_2, \ldots, x_{n-1}) \in \mathbb{R}^{n-1} \mid (x', 0) \in W \}$.
Since $F(y_1, \ldots, y_n) = (y_1, \ldots, y_{n-1}, f(y_1, \ldots, y_n)) = (y_1, \ldots, y_{n-1}, 0) \in W$, $U$ is an open set containing $(y_1, y_2, \ldots, y_{n-1}) = y'$.

Define $g: U \to \mathbb{R}$ by $g(x') = g(x_1, \ldots, x_{n-1}) = h_n(x_1, \ldots, x_{n-1}, 0)$.
Note that
$$y=\underbrace{(y',y_n)}_{\in V} = F^{-1}F(y)=F^{-1}(y',0)=(y',h_n(y',0))=(y',g(y)).$$
So $y_n = g(y').$


\vspace{0.8em}
\noindent \textbf{Step 4: Show the zero set is locally the graph of $g$}\\
\textbf{Claim:} $\{ x \in V : f(x) = 0 \} = \{ (x', g(x')) : x' \in U \}$.
Namely, we want to show the solution set of the equation $f(x)=0\,(A)$ is exactly the graph of the function $g(x')\,(B)$.
Recall $U=\{x'\in\mathbb{R}^{n-1}: (x',0)\in W\}.$
\begin{itemize}
    \item[$(\subseteq)$] Take $x = (x', x_n) \in A$, where $x\in V$ and $f(x) = 0$.\\
    Apply $F$ first, $F(x)=(x',f(x))=(x',0),$ where $F(x)\in W$, so $x'\in U.$\\
    Apply $F^{-1}$ sequentially,
    $\underbrace{F^{-1}F(x)}_{x}=F^{-1}(x',0)=(x',h_n(x',0))=(x',g(x')).$
    So $x\in B.$
    \item[$(\supseteq)$] Take $x = (x', x_n) \in B$, where $x'\in U$ and $x = (x',g(x'))$.\\
    Since $x'\in U$, $(x',0)\in W$. Apply $F^{-1}$ first, $F^{-1}(x',0)=(x',h_n(x',0))=(x',g(x'))=x,$ where $x\in V$.\\
    Apply $F$ sequentially,
    $\underbrace{FF^{-1}(x',0)}_{(x',0)}=F(x)=(x',f(x)).$
    So $x\in A.$
\end{itemize}

\vspace{0.8em}
\noindent \textbf{Step 5: Obtain IFT formula}\\
Now we have $(x_1, \dots, x_{n-1}) \in U \subseteq \mathbb{R}^{n-1}$ (where $U$ is open), $f(x', g(x')) = 0$, and $y_n = g(y')$.

Since $U$ is open, there exists $\delta > 0$ such that if $|t| < \delta$, then $y' + t e_j \in U$, which implies:
$$ f(y' + t e_j, g(y' + t e_j)) = 0 $$
Let $\phi(t) = f(y' + t e_j, g(y' + t e_j))$. Thus, $\phi(t) = 0$ when $|t| < \delta$.
Therefore, $\phi'(t) = 0$ for $|t| < \delta$.

Differentiating with respect to $t$ on both sides by the chain rule:
$$ \phi'(t) = \frac{\partial f}{\partial x_j}(y) \frac{dx_j}{dt}(t) + \frac{\partial f}{\partial x_n}(y) \cdot \frac{\partial g}{\partial x_j}(y') \frac{dx_j}{dt}(t).$$
Since $\frac{dx_j}{dt} = 1$, we evaluate at $t=0$:
$$ 0 = \frac{\partial f}{\partial x_j}(y) + \underbrace{\frac{\partial f}{\partial x_n}(y)}_{\neq 0} \frac{\partial g}{\partial x_j}(y') $$
$$ \therefore \frac{\partial g}{\partial x_j}(y') = - \frac{\frac{\partial f}{\partial x_j}(y)}{\frac{\partial f}{\partial x_n}(y)} $$


\end{pfbox}



\begin{definition}[Manifold]
A subset $M \subseteq \mathbb{R}^n$ is called a $k$-dimensional manifold if, for every point $p \in M$, there exists a neighborhood $V$ of $p$ in $M$ that is \textbf{diffeomorphic} to an open set $U$ in $\mathbb{R}^k$.
\end{definition}

\begin{remark}
In advanced calculus (e.g., Spivak's \textit{Calculus on Manifolds}), being a manifold is often established through equivalent conditions. For a subset $M \subseteq \mathbb{R}^n$ to be a $k$-dimensional manifold locally around a point $p$, the following three perspectives are mathematically equivalent:

\begin{enumerate}
    \item \textbf{The Geometric View (Diffeomorphism):} Locally, $M$ is diffeomorphic to an open set in $\mathbb{R}^k$.
    \item \textbf{The Analytical View (Graph of a Function):} Locally, $M$ can be written as the \textbf{graph} of a $C^1$ function (by expressing $n-k$ variables in terms of the other $k$ variables).
    \item \textbf{The Algebraic View (Level Set / Zero Set):} Locally, $M$ is the solution set of $n-k$ equations where the Jacobian matrix has full rank.
\end{enumerate}
\end{remark}

\subsubsection*{Applications in Economics}
\begin{enumerate}
\item Indifference Curves and Utility Functions\\
Let $U(x, y)$ be a utility function representing a consumer's preferences over two goods. An \textbf{indifference curve} is fundamentally a 1D \textbf{level curve} defined by the equation $U(x, y) = c$, where $c$ is a constant utility level. 
By the IFT, provided the marginal utility $\frac{\partial U}{\partial y} \neq 0$, we can locally express $y$ as a smooth function of $x$ (i.e., $y = g(x)$). Most importantly, the theorem gives us the exact slope of this curve, which economists call the \textbf{Marginal Rate of Substitution (MRS)}:
$$ \frac{dy}{dx} = -\frac{\frac{\partial U}{\partial x}}{\frac{\partial U}{\partial y}} = -\frac{MU_x}{MU_y} $$
The IFT guarantees we can compute this trade-off without ever needing to explicitly isolate $y$.

\item Isoquant Curves and Production Functions\\
Similarly, in producer theory, a production function is given by $Q = F(K, L)$, where $K$ is capital and $L$ is labor. An \textbf{isoquant} is the level curve $F(K, L) = c$. 
Applying the IFT, we can find the slope of the isoquant, known as the \textbf{Marginal Rate of Technical Substitution (MRTS)}:
$$ \frac{dK}{dL} = -\frac{\frac{\partial F}{\partial L}}{\frac{\partial F}{\partial K}} = -\frac{MP_L}{MP_K} $$
This tells the firm exactly how much capital can be replaced by one unit of labor while maintaining the same output level $c$.

\item Comparative Statics\\
Beyond 2D level curves, the generalized $n$-dimensional IFT is the foundation of \textbf{Comparative Statics} in macroeconomics and general equilibrium theory. 
Suppose an economic equilibrium is defined by a complex system of equations $F(\mathbf{y}, \mathbf{x}) = \mathbf{0}$, where $\mathbf{y}$ are endogenous variables (e.g., prices, GDP) and $\mathbf{x}$ are exogenous policy parameters (e.g., tax rates, money supply). Even if the system is too complex to explicitly solve for the equilibrium state $\mathbf{y}^* = g(\mathbf{x})$, the IFT guarantees that this relationship exists locally, provided the Jacobian matrix $D_{\mathbf{y}}F$ is invertible. Furthermore, it allows economists to calculate $\frac{\partial \mathbf{y}^*}{\partial \mathbf{x}}$, precisely predicting how a shift in government policy will ripple through the equilibrium variables.
\end{enumerate}